\chapter{Conclusion}
\label{ch:conclusion}

So what did we actually build, and does it work?

At its core, this thesis introduced a closed-loop Retrieval-Augmented Generation framework—one that does not just \emph{detect} hallucination but actually \emph{does something about it} before a risky response ever reaches the user. The system instruments frozen decoder-only language models with PyTorch forward hooks, extracts two complementary activation signals (CEV from the residual stream, IAV from the SwiGLU feed-forward intermediate), trains lightweight MLP probes on those signals, fuses their calibrated outputs into a single hallucination-risk estimate, and feeds that estimate into a deterministic controller. The controller can accept, regenerate, re-retrieve, or abstain—turning passive detection into active inference-time control.

The empirical results are encouraging. Across three open-weight decoder backbones, the fused CEV/IAV probe achieves strong in-domain validation performance: \textbf{0.8596 AUROC} on Mistral-7B-Instruct-v0.2, \textbf{0.8808} on Qwen3-8B, and \textbf{0.8689} on LLaMA-3-8B-Instruct. These numbers surpass every published and reproduced baseline we compared against—and ours is the only method in the comparison that actually closes the loop.

More importantly, the closed-loop evaluation shows that detection can translate into risk reduction. The controller cuts delivered hallucination proxy by 45--79\% depending on the backbone, with each model landing at a different point on the safety--utility curve (Table~\ref{tab:cross_model}, Section~\ref{sec:safety_utility}). Mistral balances risk and throughput. Qwen maximizes safety at the cost of frequent refusals. LLaMA maximizes utility but accepts more residual risk.

The HaluEval analysis taught us something sobering about generalization. LLaMA and Qwen transfer their learned signals cleanly to out-of-domain data. Mistral's signal is equally strong but polarity-inverted—a fascinating result that underscores the need for target-domain validation before any real deployment.

\textbf{Where should future work go?} Several directions seem most pressing: evaluation on the official RAGTruth test split rather than just a validation partition; multi-seed replication to tighten confidence intervals; stronger target-domain calibration for the Mistral polarity issue; content-quality checks that go beyond hallucination to measure omission, vagueness, and relevance failures; and extension to non-SwiGLU architectures to test broader generality.

%% -------------------------------------------------
\section{Architecture-Agnostic Scope}
\label{sec:agnostic_scope}

One question that naturally arises: is this framework tied to a specific model, or can it be reused?

The answer is both. The \emph{design} is architecture-agnostic—the same five-step procedure applies regardless of which backbone you plug in:

\begin{enumerate}
    \item Register hooks at proportional decoder depths.
    \item Extract CEV and IAV features.
    \item Train lightweight MLP probes on hallucination labels.
    \item Calibrate and fuse the probe probabilities.
    \item Feed the fused score into the same four-action controller.
\end{enumerate}

\noindent But the \emph{weights} are not transferable. Each model has its own activation geometry, its own scale, its own tokenizer quirks, and its own FFN intermediate dimensions. You cannot train probes on Mistral and deploy them on Qwen—the activations live in different spaces. So the framework is reusable in the sense that the recipe works everywhere, but each new backbone requires its own round of probe training and threshold calibration.

Table~\ref{tab:agnostic} makes this explicit.

\begin{table}[!htb]
\centering
\caption{What stays fixed and what changes when you swap backbones.}
\label{tab:agnostic}
\begin{tabular}{@{}lcp{7cm}@{}}
\toprule
\textbf{Component} & \textbf{Shared?} & \textbf{Explanation} \\
\midrule
Hooking strategy    & Yes   & Same proportional-depth approach ($L/4$, $L/2$, $3L/4$) \\
CEV/IAV concept     & Yes   & Same signals—residual stream and FFN intermediate \\
Probe architecture  & Mostly & Same MLP structure; input dimension adapts to backbone \\
Controller logic    & Yes   & Same four-action deterministic policy \\
Threshold values    & No    & Must be calibrated per backbone \\
Probe weights       & No    & Must be trained from scratch for each model \\
Input dimensions    & No    & Depend on $d_{\text{hidden}}$ and $d_{\text{inter}}$ \\
\bottomrule
\end{tabular}
\end{table}

\textbf{In plain terms:} the system is architecturally portable across related decoder-only transformer families. But portability does not mean zero-cost transfer—each deployment needs its own calibration pass. That is an honest characterization of where internal-state probing stands today, and it is a limitation worth highlighting rather than glossing over.
